% Generated draft: requires mathematical and source review.
\documentclass[11pt]{article}
\usepackage[margin=25mm]{geometry}
\usepackage{fontspec}
\setmainfont{Latin Modern Roman}
\usepackage{amsmath,amssymb,mathtools,unicode-math}
\setmathfont{Latin Modern Math}
\usepackage{enumitem,needspace,xurl,hyperref}
\setlength{\parindent}{0pt}
\setlength{\parskip}{5pt}
\setlength{\emergencystretch}{3em}
\newcommand{\N}{\mathbb{N}}
\newcommand{\Z}{\mathbb{Z}}
\newcommand{\Q}{\mathbb{Q}}
\newcommand{\R}{\mathbb{R}}
\newcommand{\C}{\mathbb{C}}
\newcommand{\F}{\mathbb{F}}
\newcommand{\A}{\mathbb{A}}
\newcommand{\PP}{\mathbb{P}}
\newcommand{\E}{\mathbb{E}}
\newcommand{\eps}{\varepsilon}
\newcommand{\dd}{\,\mathrm{d}}
\begin{document}
\section*{UnsolvedMath ID: 30002333}
\textbf{Source problem code:} catalogue entry
\subsection{Definitions and mathematical statement}
Let \([L]\) denote a two-component link type, and let \(\mathcal{R}([L])\) denote its collection of admissible geometric representatives \(L=K_1\sqcup K_2\), with admissibility and equivalence understood as in the cited source. The supplied record calls \(A(K_1,K_2)\) the \emph{symplectic measure} of a representative but does not define this functional. Formally, the link-type quantity under consideration is
\[
A([L]) := \inf\bigl\{A(K_1,K_2):K_1\sqcup K_2\in\mathcal{R}([L])\bigr\}.
\]
The requested problem is to determine \(A([L])\) for two-component link types and, in particular, to decide whether
\[
\forall\,[L]\text{ a hyperbolic two-component link type},\qquad A([L])>0.
\]
Here the infimum is over representatives of each fixed link type; the positivity question is posed separately for every hyperbolic type.

\textbf{Unresolved qualifications.} This is a formal transcription of the intended question, not a complete definition of its invariant. The supplied record does not specify the formula or normalization for \(A(K_1,K_2)\), its domain and range, the ambient space, the regularity of representatives, or the equivalence relation defining link type. It also does not define the intended hyperbolicity condition. These source-dependent choices must be recovered before the displayed question can be treated as a fully specified mathematical statement.
\subsection{Short English statement}
For each two-component link type, take the infimum of the symplectic measure over all its admissible representatives. Determine this infimum and ask whether it is strictly positive for every hyperbolic two-component link type. The supplied record omits the definition of the measure and the precise admissibility and hyperbolicity conventions, so the statement requires source clarification.
\subsection{Sources}
\begin{itemize}
\item Supplied local UnsolvedMath registry record, permanent ID \(30002333\), source problem code \texttt{catalogue entry}, titled \emph{Symplectic Measure of Two-Component Link Types}, in the Oberwolfach Reports Open Problems collection. The record gives proposed year 2013 and names no proposer.
\item Supplied source citation: \emph{Geometric Knot Theory (2014)}, Oberwolfach Reports, pp.\ 1313--1358, DOI \texttt{10.4171/owr/2013/22}. The literature-source list also calls this \emph{Oberwolfach Report 22/2013}. \url{https://doi.org/10.4171/owr/2013/22}. The report is described as documenting workshop discussions and open problems.
\item Supplied literature reference: O'Hara, \emph{Measure of a 2-component link}. \url{https://doi.org/10.2748/tmj/1378991024}.
\item Supplied literature assessment, checked on \textbf{2026-08-21}: the record labels the problem open and attributes to O'Hara a characterization of the zero minimum for a fixed geometric link and known zero-infimum classes, while describing positivity of \(A([L])\) for hyperbolic links as open. This is supplied dated curation derived from the OWR record, not independently verified literature assessment. The cited claims and time-sensitive status require independent verification.
\end{itemize}
\end{document}
